\subsection{Area amministratore}

Se accedi al sistema con un account \hyperlink{glossario:amministratore}{amministratore}, vedrai un'\hyperlink{glossario:interfaccia}{interfaccia} diversa rispetto a quella dell'utente standard.\\
L'area \hyperlink{glossario:amministratore}{amministratore} \`e pensata per monitorare tutti gli ordini del sistema, senza utilizzare la \hyperlink{glossario:chat}{chat} o il carrello.

\subsubsection{Accesso amministratore}

Accedendo con credenziali da \hyperlink{glossario:amministratore}{amministratore}, verrai portato direttamente alla schermata dello \hyperlink{glossario:storico-ordini}{storico ordini} globale, che rappresenta l'area principale di lavoro.

% TODO: inserire figura schermata iniziale area amministratore

A differenza dell'\hyperlink{glossario:utente-standard}{utente standard}:
\begin{itemize}[itemsep=5pt, parsep=5pt, label=$\scriptstyle\bullet$]
	\item Non \`e presente la \hyperlink{glossario:chat}{chat}
	\item Non \`e presente il \hyperlink{glossario:carrello}{carrello}
	\item Non \`e possibile creare o modificare ordini
\end{itemize}

L'\hyperlink{glossario:amministratore}{amministratore} utilizza quindi il sistema solo per consultare e analizzare gli ordini.

\subsubsection{Storico ordini globale}

Nella schermata principale puoi visualizzare tutti gli ordini effettuati dagli utenti.

% TODO: inserire figura tabella storico ordini globale

Ogni \hyperlink{glossario:ordine}{ordine} \`e mostrato in una tabella con le informazioni principali:
\begin{itemize}[itemsep=5pt, parsep=5pt, label=$\scriptstyle\bullet$]
	\item Codice \hyperlink{glossario:ordine}{ordine}
	\item Data
	\item Cliente
	\item Numero di prodotti
	\item Accesso al dettaglio dell'\hyperlink{glossario:ordine}{ordine}
\end{itemize}

La presenza della colonna cliente ti permette di capire facilmente chi ha effettuato ciascun \hyperlink{glossario:ordine}{ordine}, cosa che non \`e visibile nello storico dell'utente standard.

\subsubsection{Filtri}

Per trovare pi\`u facilmente gli ordini, puoi utilizzare il filtro per data.\\
Ti basta selezionare un intervallo temporale per visualizzare solo gli ordini effettuati in quel periodo.

% TODO: inserire figura filtro per data area amministratore

Questo ti permette di concentrarti su un arco di tempo specifico e semplifica la ricerca delle informazioni.

\subsubsection{Paginazione}

Se il numero di ordini \`e elevato, i risultati vengono suddivisi in pi\`u pagine.\\
Nella parte inferiore della schermata trovi i pulsanti:
\begin{itemize}[itemsep=5pt, parsep=5pt, label=$\scriptstyle\bullet$]
	\item \vr{Precedente}
	\item \vr{Successivo}
\end{itemize}

che ti permettono di scorrere tra le diverse pagine dello storico.

% TODO: inserire figura paginazione area amministratore

\subsubsection{Limitazioni amministratore}

Come \hyperlink{glossario:amministratore}{amministratore}, utilizzi il sistema solo per consultare gli ordini.\\
Per questo motivo alcune funzionalit\`a non sono disponibili:

\begin{itemize}[itemsep=5pt, parsep=5pt, label=$\scriptstyle\bullet$]
	\item \textbf{Nessuna \hyperlink{glossario:chat}{chat}}: non puoi interagire con l'assistente AI.
	\item \textbf{Nessun \hyperlink{glossario:carrello}{carrello}}: non \`e presente la sezione per gestire prodotti.
	\item \textbf{Nessuna creazione di ordini}: non puoi aggiungere prodotti o inviare nuovi ordini.
\end{itemize}

Questa separazione permette di distinguere chiaramente tra chi utilizza il sistema per creare ordini e chi lo utilizza per monitorarli.
